% ======================================================================
% SECTION 1: INTRODUCTION
% Five thematic subsections expanded from the v0.8.0 bracketed draft at
% 2030-pdac-1min/paper/draft-paper/sections/introduction.tex.
% Location: 2030-pdac-1min/paper/full-paper/sections/introduction.tex
% ======================================================================

\subsection{FDA Real-Time Clinical Trials Announcement and Surgical Extension}
\label{sec:intro-fda}

On 28 April 2026 the United States Food and Drug Administration
(FDA) announced that real time clinical trials (RTCT) are entering
the proof of concept phase \cite{fda2026realtime}. The announcement
names two oncology pharmacology pilots, AstraZeneca TRAVERSE in
mantle cell lymphoma and Amgen STREAM SCLC in small cell lung
carcinoma.\ FDA Commissioner Marty Makary framed the program as a
structural upgrade to the clinical trial data pipeline, and FDA
Chief AI Officer Jeremy Walsh tied the program to the agency
software and AI modernization roadmap.\ The prior 60 second
glioblastoma resection paper \cite{kawchak_2026_20113157} positioned
that announcement as a natural inflection for the surgical theater
through an on premises repository based LLM thesis. The present
work continues that opportunistic surgical extension by pairing a
precision oncology drug development program (RASolute 302 plus
RASolve 301 \cite{rasolute302, rasolve301}) with a precision
oncology robotic surgery program (hypothetical 2030 Medtronic
PancreSpeed 1.0).

The work is independent and is not endorsed or sponsored by trial
sponsors, the FDA, contract research organizations, sites,
institutional review boards, regulators, or medical societies.\ The
FDA Software as a Medical Device (SaMD) framework \cite{fda-samd}
provides the regulatory anchor without claiming endorsement, and
the United States Number 1 patient safety, efficacy, and speed
framing carried over from the v0.4.0 GBM paper is preserved
verbatim in this paper. The aim of this work is to publish a
reproducible computational checkpoint that the industry can pick
up rather than re author from scratch, in the same spirit that
the GBM paper served the neurosurgical community.

The remainder of this section frames the clinical context for
pancreatic ductal adenocarcinoma (\S \ref{sec:intro-pdac}), the
current robotic Whipple state of the art (\S
\ref{sec:intro-baseline}), the on premises repository based LLM
thesis (\S \ref{sec:intro-thesis}), and the transition to the 60
second 8 arm Whipple plus Daraxonrasib pairing (\S
\ref{sec:intro-transition}).

\subsection{Pancreatic Cancer Clinical Context}
\label{sec:intro-pdac}

Pancreatic ductal adenocarcinoma (PDAC) is the third leading cause
of cancer death in the United States and the fourth in the
European Union, with a total 5 year overall survival below 13 percent
\cite{Siegel2025CancerStatistics}. The pancreaticoduodenectomy
(Whipple procedure) is the most technically demanding of all
curative intent abdominal oncology operations. It requires
resection plus reconstruction of four luminal structures (the
pancreas, the duodenum, the common bile duct, and the stomach or
pylorus) and dissection of two named vessels (the superior
mesenteric vein and the portal vein) plus the artery first
approach to the superior mesenteric artery. The leading lethal
postoperative cascade is pancreatic fistula, intra abdominal
infection, hemorrhage, sepsis, and failure to rescue, scored under
the International Study Group on Pancreatic Surgery (ISGPS) Grade
A, B, and C postoperative pancreatic fistula classification
\cite{Bassi2017ISGPSPostOpFistula}.

The 2025 Dutch nationwide cohort of 1000 robotic
pancreaticoduodenectomies reported conversion of 10.1 percent,
Clavien Dindo grade III or higher complications of 41.3 percent,
ISGPS grade B or C postoperative pancreatic fistula of 24.4
percent, 30 day mortality of 3.9 percent, and a mean ideal outcome
rate of 47 percent \cite{DutchCohort2025Whipple}. These numbers
anchor the human baseline used by every comparison in this paper.

The synthetic patient PAT-PDAC-0001 used throughout this work is a
68 year old with a 3.4 cm head of pancreas PDAC abutting the
superior mesenteric vein at 75 degrees, KRAS G12D mutant, MSI
stable, CA 19-9 of 412 U/mL at diagnosis, ECOG performance status
1, and a completed neoadjuvant modified FOLFIRINOX course
(times 4 cycles) \cite{Conroy2018FOLFIRINOXAdjuvant}. The patient
is Daraxonrasib eligible under the broad pan KRAS criteria of
RASolute 302 \cite{rasolute302}. The identity, disease parameters,
and pre operative imaging are inherited verbatim from
\texttt{2030-pdac-1min/paper/instructions/pdac_context_1min.md}.
Current operative duration for a single arm robotic Whipple ranges
between 4 and 8 hours of in room time; a 60 second target is a
240x to 480x compression that no current commercial robot
achieves.\ The contrast frames the rest of the paper as a checking
step the industry has not yet taken.

\clearpage

\subsection{Current Robotic Whipple State of the Art and the 4 Entrant Tournament}
\label{sec:intro-baseline}

The state of the art for abdominal robotic surgery in 2026 spans
three commercial platforms cited in
\texttt{2030-pdac-1min/paper/instructions/competition_protocol.md}:
the Intuitive da Vinci SP single port system
\cite{intuitive-davinci-sp}, and the Medtronic Hugo RAS multi arm
platform \cite{medtronic-hugo-ras}.\ None can perform a 60
second Whipple because per arm end effector velocity, joint
angular velocity, E stop latency, vessel surface positioning
accuracy, and force resolution are each 5x to 500x short of the
requirement. The hypothetical 2030 PancreSpeed 1.0 closes the
gaps by design: 8 cooperating arms, 1,200 mm/s peak tip velocity,
100 kHz force sampling, 1,600 mm cubed per second peak removal,
0.05 mm RMS positioning at 1,200 mm/s, 3 ms E stop, an 18 N
cumulative cross arm tip force cap, and a 3 N per arm tip force
cap.

\begin{table}[H]
\caption{Robot platform comparison: 2025 state of the art versus the 2030 PancreSpeed 1.0 target.}
\label{tab:robot-comparison}
\centering
\begin{tabular}{>{\raggedright\arraybackslash}p{3.7cm} >{\raggedright\arraybackslash}p{4.7cm} >{\raggedright\arraybackslash}p{6.8cm}}
\toprule
\textbf{Spec} & \textbf{2025 SOTA (da Vinci SP, Hugo RAS)} & \textbf{2030 PancreSpeed 1.0 target} \\
\midrule
  Arms & 1 to 4 \cite{intuitive-davinci-sp, medtronic-hugo-ras} & 8 (dissection, three anastomoses, hemostasis, drains) \\
  Tip velocity & up to 250 mm/s & 1,200 mm/s (about 5x faster) \\
  Force sampling & 1 to 10 kHz & 100 kHz force plus 10 kHz command (10x finer) \\
  E stop latency & 30 to 50 ms & 3 ms (10x to 17x faster) \\
  Position accuracy at velocity & 0.5 to 1.0 mm RMS & 0.05 mm RMS at 1,200 mm/s (10x to 20x finer) \\
\bottomrule
\end{tabular}
\end{table}

The 4 entrant tournament protocol defined in
\texttt{2030-pdac-1min/paper/instructions/competition_protocol.md}
and the per round verdict format in
\texttt{commit_05_competition_1min.md} score PancreSpeed 1.0
against three hypothetical 2030 successor robots (da Vinci SP
successor, Hugo RAS oncology successor)
and the 2025 Dutch human surgeon baseline
\cite{DutchCohort2025Whipple}. The structural time dimension
caveat (1 minute robot versus 4 to 8 hour human baseline) is
preserved in every per round LLM judge rationale. The next
subsection states the on premises repository based LLM thesis
that the tournament tests.

\subsection{On-Premises Repository Based LLM Thesis}
\label{sec:intro-thesis}

The thesis is stated verbatim from
\texttt{2030-pdac-1min/paper/instructions/README.md}: on premises
repository based large language models provide commands to
standard oncology surgical robots based on real time sensor data
and are controlled via x, y, z coordinates to administer patient
treatment, and this workflow minimizes single robot error
potential.\ The mechanism is second opinion oracle isolation.\ The
LLM sits outside the robot vendor kinematic stack, so a fault in
the vendor controller does not propagate into the LLM advisory
layer.\ Cross referenced with
\texttt{2030-pdac-1min/paper/instructions/multi_arm_coordination_8arm.md},
LLM advisories are required to respect a 10 kHz heartbeat
broadcast bus, a 100 microsecond watchdog, and a 3 ms cross arm E
stop budget.

\clearpage

\begin{verbatim}
+-------------------------------------------------------------------+
|  10 kHz heartbeat broadcast bus (64 byte frame, 100 us deadline)  |
+-------------------------------------------------------------------+
|  Arm 1 SMV dissect |  Arm 2 PV dissect  |  Arm 3 hep artery ctrl  |
|  hybrid u-w-p tip  |  hybrid u-w-p tip  |  bipolar + retractor    |
|  <= 3 N per arm    |  <= 3 N per arm    |  <= 3 N per arm         |
+-------------------------------------------------------------------+
|  Arm 4 NIR + ICG   |  Arm 5 PJ ring     |  Arm 6 bipolar coag     |
|  imaging probe     |  tension monitor   |  + cautery              |
|  <= 0.5 N contact  |  <= 2.5 N anast    |  <= 3 N per arm         |
+-------------------------------------------------------------------+
|  Arm 7 suction     |  Arm 8 final image + drain placement         |
|  + irrigation      |  + closing NIR scan                          |
+-------------------------------------------------------------------+
|  Cumulative cross arm tip force on patient frame <= 18 N          |
|  E stop 3 ms / heartbeat watchdog 100 us / emergency arm park     |
+-------------------------------------------------------------------+
\end{verbatim}
 

Every LLM command is sourced from a hash pinned in GitHub repository files under locations
\texttt{2030-pdac-1min/paper/codegen/} or
\texttt{2030-pdac-1min/paper/execution/}, so each command can be
re traced to a deterministic seed and to a specific commit.
On premises inference backends carry the same Claude Code agent stack
\cite{claude-code, claude-opus-47, claude-sonnet-46} used to
generate the five phase workflow in this paper, while preserving
the data residency requirement that aligns with FDA SaMD
\cite{fda-samd}, ICH E6(R3) Good Clinical Practice
\cite{ich-e6r3}, and 21 CFR 50.30 informed consent
\cite{cfr-21-50-30}.

\subsection{Transition to the 60-Second 8-Arm Whipple Plus Daraxonrasib Pairing}
\label{sec:intro-transition}

The artifacts paired with this paper are the output of a five
phase Claude Code workflow. Phase 1 is the v0.5.0 instruction set
at \texttt{2030-pdac-1min/paper/instructions/}
\cite{pdac-instructions-v050}, phase 2 is the v0.6.0 codegen tree
at \texttt{2030-pdac-1min/paper/codegen/}
\cite{pdac-codegen-v060}, phase 3 is the v0.7.0 execution tree at
\texttt{2030-pdac-1min/paper/execution/}
\cite{pdac-execution-v070}, and phase 4 is the v0.8.0 bracketed
draft template at \texttt{2030-pdac-1min/paper/draft-paper/}
\cite{pdac-draft-paper-v080} that this v0.9.0 full paper expands
without modifying the upstream draft.

The Daraxonrasib rationale is anchored in
\texttt{2030-pdac-1min/paper/instructions/daraxonrasib_integration.md} 
and the Daraxonrasib medicine large language model trial summary
\cite{kawchak_2025_18099351} at
\texttt{2030-pdac-1min/paper/inputs/daraxonrasib-1/main.tex}.
PAT-PDAC-0001 is KRAS G12D positive and Daraxonrasib eligible per
the RASolute 302 broad pan KRAS criteria \cite{rasolute302}; the
RASolve 301 first line program \cite{rasolve301} adds a checkpoint
inhibitor combination that the perioperative pause and restart
logic respects. FDA Breakthrough Therapy designation in June 2025
\cite{rev-fda-breakthrough} sets the regulatory anchor without
endorsement. The on premises LLM bound advisory layer recommends a
postoperative restart day on T+7d, T+14d, or T+21d based on per
iteration sensor and anastomosis outcomes (see \S
\ref{sec:methods-daraxonrasib}). The four prior author PDAC papers
under \texttt{2030-pdac-1min/paper/inputs/paper-1/} through
\texttt{paper-4/} \cite{kawchak_2025_15735068,
kawchak_2025_16415815, kawchak_2025_17001137,
kawchak_2025_17239510} independently complemented real RASolute
302 trial development throughout 2025; the present paper makes the
surgical robot complement explicit and pairs the pharmacologic and
surgical workflows under one repository. The Methods section
details the deterministic seed contract (root seed 20260513), the
32 iteration sweep, and the 4 entrant tournament protocol.

\clearpage
